\documentclass[10pt,a4paper]{amsart}
\usepackage{amsmath,amssymb,amsthm,graphicx}
% For commutative diagrams we use Xy-pic
\usepackage[all]{xy}

% We use 2cell for 2-commutative diagrams.
\xyoption{2cell}
% Theorem environments.
%
\theoremstyle{plain}
\newtheorem{theorem}[subsection]{Theorem}
\newtheorem{proposition}[subsection]{Proposition}
\newtheorem{lemma}[subsection]{Lemma}

\theoremstyle{definition}
\newtheorem{definition}[subsection]{Definition}
\newtheorem{example}[subsection]{Example}
\newtheorem{exercise}[subsection]{Exercise}
\newtheorem{situation}[subsection]{Situation}

\theoremstyle{remark}
\newtheorem{remark}[subsection]{Remark}
\newtheorem{remarks}[subsection]{Remarks}

\numberwithin{equation}{subsection}

% Macros
%
\def\lim{\mathop{\mathrm{lim}}\nolimits}
\def\colim{\mathop{\mathrm{colim}}\nolimits}
\def\Spec{\mathop{\mathrm{Spec}}}
\def\Hom{\mathop{\mathrm{Hom}}\nolimits}
\def\Ext{\mathop{\mathrm{Ext}}\nolimits}
\def\SheafHom{\mathop{\mathcal{H}\!\mathit{om}}\nolimits}
\def\SheafExt{\mathop{\mathcal{E}\!\mathit{xt}}\nolimits}
\def\Sch{\mathit{Sch}}
\def\Mor{\mathop{\mathrm{Mor}}\nolimits}
\def\Ob{\mathop{\mathrm{Ob}}\nolimits}
\def\Sh{\mathop{\mathit{Sh}}\nolimits}
\def\NL{\mathop{N\!L}\nolimits}
\def\CH{\mathop{\mathrm{CH}}\nolimits}
\def\proetale{{pro\text{-}\acute{e}tale}}
\def\etale{{\acute{e}tale}}
\def\QCoh{\mathit{QCoh}}
\def\Ker{\mathop{\mathrm{Ker}}}
\def\Im{\mathop{\mathrm{Im}}}
\def\Coker{\mathop{\mathrm{Coker}}}
\def\Coim{\mathop{\mathrm{Coim}}}

% Boxtimes
%
\DeclareMathSymbol{\boxtimes}{\mathbin}{AMSa}{"02}

%
% Macros for moduli stacks/spaces
%
\def\QCohstack{\mathcal{QC}\!\mathit{oh}}
\def\Cohstack{\mathcal{C}\!\mathit{oh}}
\def\Spacesstack{\mathcal{S}\!\mathit{paces}}
\def\Quotfunctor{\mathrm{Quot}}
\def\Hilbfunctor{\mathrm{Hilb}}
\def\Curvesstack{\mathcal{C}\!\mathit{urves}}
\def\Polarizedstack{\mathcal{P}\!\mathit{olarized}}
\def\Complexesstack{\mathcal{C}\!\mathit{omplexes}}
% \Pic is the operator that assigns to X its picard group, usage \Pic(X)
% \Picardstack_{X/B} denotes the Picard stack of X over B
% \Picardfunctor_{X/B} denotes the Picard functor of X over B
\def\Pic{\mathop{\mathrm{Pic}}\nolimits}
\def\Picardstack{\mathcal{P}\!\mathit{ic}}
\def\Picardfunctor{\mathrm{Pic}}
\def\Deformationcategory{\mathcal{D}\!\mathit{ef}}
\title{The determinant module of a finite length module over a local ring is at most one dimensional over the residue field}
\author{The Stacks Project Authors}
\date{stacks-project:chow.tex:local-snapshot-20261012; GNU FDL 1.2 or later}
\begin{document}
\maketitle
\noindent\textit{Editorial note.} All delivered mathematical text is the author's own, copied from the pinned Stacks Project chapter file chow.tex; no statement or proof step was rewritten, reordered or removed.\par\medskip
\noindent\textit{Editorial note.} Preamble provenance: the pinned chapter file opens with an input of the shared preamble and carries no preamble text of its own, so the project's single Stacks-wide preamble staged at sources/stacks-chow/preamble.tex is excerpted verbatim as two spans, lines 14-18 (Xy-pic setup) and lines 154-217 (the native theorem declarations and mathematical macros). Nothing was added, removed or redefined; the wrapper class and the amsmath/amssymb/amsthm/graphicx base packages are supplied by the assembler.\par\medskip
\noindent\textit{Editorial note.} The counted statement and proof contain no cross reference, no citation and no included file, so the delivered body is self contained.\par\medskip
\noindent\textit{Editorial note.} Canonical location of the counted proposition: Stacks tag 02P7 (sources/stacks-chow/pinned-original-source.tex line 15761).\par\medskip
\noindent\textit{Allowed standard prerequisites (editorial).}\begin{itemize}
\item The determinant module det\_kappa(M) of a finite length module over a local ring, with admissible sequences and the three admissible relations. Same chapter, chow.tex, Section 'Periodic complexes' and the definitions of admissible sequences and of det\_kappa preceding this lemma (the statement's phrase 'with notations as above' refers to them). Supplies the notion of admissible sequence, the three relations (vanishing when the sequence does not generate M, scaling by a unit and moving an element that is killed by m), and the fact that det\_kappa(M) is generated by the symbols of admissible sequences; the replacement argument itself is carried out in the counted span.
\end{itemize}\medskip
\begin{lemma}
\label{lemma-dimension-at-most-one}
With notations as above we have $\dim_\kappa(\det_\kappa(M)) \leq 1$.
\end{lemma}

\begin{proof}
Fix an admissible sequence $(f_1, \ldots, f_l)$ of $M$ such that
$$
\text{length}_R(\langle f_1, \ldots, f_i\rangle) = i
$$
for $i = 1, \ldots, l$. Such an admissible sequence exists exactly because
$M$ has length $l$. We will show that any element of
$\det_\kappa(M)$ is a $\kappa$-multiple of the symbol
$[f_1, \ldots, f_l]$. This will prove the lemma.

\medskip\noindent
Let $(e_1, \ldots, e_l)$ be an admissible sequence of $M$.
It suffices to show that $[e_1, \ldots, e_l]$ is a multiple
of $[f_1, \ldots, f_l]$. First assume that
$\langle e_1, \ldots, e_l\rangle \not = M$. Then there exists
an $i \in [1, \ldots, l]$ such that
$e_i \in \langle e_1, \ldots, e_{i - 1}\rangle$. It immediately
follows from the first admissible relation that
$[e_1, \ldots, e_n] = 0$ in $\det_\kappa(M)$.
Hence we may assume that $\langle e_1, \ldots, e_l\rangle = M$.
In particular there exists a smallest index $i \in \{1, \ldots, l\}$
such that $f_1 \in \langle e_1, \ldots, e_i\rangle$. This means
that $e_i = \lambda f_1 + x$ with
$x \in \langle e_1, \ldots, e_{i - 1}\rangle$ and $\lambda \in R^*$.
By the second admissible relation this means that
$[e_1, \ldots, e_l] =
\overline{\lambda}[e_1, \ldots, e_{i - 1}, f_1, e_{i + 1}, \ldots, e_l]$.
Note that $\mathfrak m f_1 = 0$. Hence by applying the third
admissible relation $i - 1$ times we see that
$$
[e_1, \ldots, e_l] =
(-1)^{i - 1}\overline{\lambda}
[f_1, e_1, \ldots, e_{i - 1}, e_{i + 1}, \ldots, e_l].
$$
Note that it is also the case that
$ \langle f_1, e_1, \ldots, e_{i - 1}, e_{i + 1}, \ldots, e_l\rangle = M$.
By induction suppose we have proven that our original
symbol is equal to a scalar times
$$
[f_1, \ldots, f_j, e_{j + 1}, \ldots, e_l]
$$
for some admissible sequence $(f_1, \ldots, f_j, e_{j + 1}, \ldots, e_l)$
whose elements generate $M$, i.e., \ with
$\langle f_1, \ldots, f_j, e_{j + 1}, \ldots, e_l\rangle = M$.
Then we find the smallest $i$ such that
$f_{j + 1} \in \langle f_1, \ldots, f_j, e_{j + 1}, \ldots, e_i\rangle$
and we go through the same process as above to see that
$$
[f_1, \ldots, f_j, e_{j + 1}, \ldots, e_l]
=
(\text{scalar}) [f_1, \ldots, f_j, f_{j + 1}, e_{j + 1},
\ldots, \hat{e_i}, \ldots, e_l]
$$
Continuing in this vein we obtain the desired result.
\end{proof}

\end{document}
